\section{The B04/A7c selected-hinge bridge}
\label{sec:b04}

The B04 rows are zero-margin selected hinge contacts.  They cannot be certified
by positive SAT clearance: the pieces are meant to remain joined along the
hinge axis.  The correct local statement is boundary contact plus no strict
local interior overlap away from the hinge axis.

Fix one selected row.  Let the oriented hinge axis be \(a\to b\), and let
\(F\) be the third vertex of the source shared face, so the source face is
\(\conv(a,b,F)\).  Let \(Q_p\) be the parent apex not in this face and \(Q_c\)
the child apex not in this face.  Define the oriented side form
\[
\Omega(X,Y)=(b-a)\cdot((X-a)\times(Y-a)).
\]

\begin{lemma}[Local selected-hinge wedge]
\label{lem:wedge}
For each of the six selected B04 rows, with A7c ray sign \(\sigma\), the package
records the exact side identities
\[
\sigma\Omega(F,Q_p)=-\frac{\sqrt2}{8}<0,
\qquad
\sigma\Omega(F,Q_c)=\frac{\sqrt2}{8}>0.
\]
Thus the A7c sign identifies the child opening side of the selected source
face.  Moreover the cross-section angle \(\alpha\) between the source face ray
and either adjacent tetrahedral apex ray satisfies
\[
\cos\alpha=\frac{\sqrt6}{3}>\frac{1}{2},
\qquad\text{so}\qquad
0<\alpha<60^\circ.
\]
Consequently, for \(0<\theta\le120^\circ\), the opened child sector satisfies
\(\theta+\alpha<180^\circ\) and cannot wrap behind the hinge axis into the
parent sector.
\end{lemma}

\begin{proof}
The identities for \(\Omega\) are exact coordinate evaluations in the normal
cross-section about the selected hinge axis.  They place the parent apex and
child apex on opposite sides of the oriented source face, after multiplying by
the row sign \(\sigma\).  The angle calculation gives
\(\cos\alpha=\sqrt6/3>1/2\), hence \(\alpha<60^\circ\).  Since the ray domain
has \(\theta\le120^\circ\), the sector opened by the child face remains within
an open half-turn from the source face.  Therefore it cannot cross through the
opposite parent sector away from the hinge axis.
\end{proof}

The bridge also checks the actual rotated child vertices using
\(t=\tan(\theta/2)\).  For the rotated third vertex of the selected source face,
the signed side is
\[
\frac{t}{2(t^2+1)}>0,
\qquad 0<t\le\sqrt3.
\]
For the rotated child apex, each row reduces to a positive rational form
recorded in the row JSON: the denominator is positive, the sample signs at
\(t=1\) and \(t=\sqrt3\) are positive, and all positive numerator roots lie
strictly above \(\sqrt3\).

\begin{certprop}[B04/A7c bridge rows]
\label{prop:b04-bridge}
The package-local bridge accepts exactly six records:
\begin{center}
\begin{tabular}{llllr}
\toprule
Tree & Pair & Hinge & Contact & Ray sign\\
\midrule
\treeA{} & \(P_0-P_1\) & H0\_A\_M\_AB & C0 & \(+1\)\\
\treeA{} & \(P_0-P_2\) & H4\_C\_M\_CD & C1 & \(+1\)\\
\treeA{} & \(P_1-P_3\) & H7\_D\_M\_CD & C4 & \(-1\)\\
\treeB{} & \(P_0-P_1\) & H0\_A\_M\_AB & C0 & \(+1\)\\
\treeB{} & \(P_1-P_3\) & H7\_D\_M\_CD & C4 & \(-1\)\\
\treeB{} & \(P_2-P_3\) & H9\_B\_M\_AB & C5 & \(+1\)\\
\bottomrule
\end{tabular}
\end{center}
For these rows, A7c signed orientation plus the local wedge bridge proves
boundary hinge contact and no strict local interior overlap at the contact
plane, excluding the hinge axis.
\end{certprop}

\begin{proof}[Certification]
The bridge manifest is
\path{certified/b04_a7c_contact_side_bridge_manifest.json}.  It records
\texttt{record\_count=6}, \texttt{accepted\_bridge\_record\_count=6}, and the
policy that bridge rows are consumed from the A7c manifest records while
asserting equality of \texttt{tree\_id}, \texttt{hinge\_id}, \texttt{pair}, and
\texttt{ray\_sign}.  The row records live in
\path{certified/a7c_b04_contact_side_bridge_records/}.  The derivation and
claim boundary are source-locked in
\path{certified/source_docs/S4_CL5_A7C_B04_CONTACT_SIDE_BRIDGE_CERTIFICATE.md}.
\end{proof}

\begin{remark}
The bridge is local and selected-row scoped.  It does not give positive
separation for selected hinges, does not handle hinge thickness, and does not
extend to unselected or three-parameter B04 cases.
\end{remark}
